% Two movements the driver makes at right angles arrive 137.7 degrees apart.
\begin{tikzpicture}[x=1mm, y=1mm, font=\scriptsize]

  % intended
  \begin{scope}[shift={(0,0)}]
    \node[font=\small\bfseries, anchor=south] at (18,26) {What the driver did};
    \draw[go, thick] (0,0) -- (24,0) node[right, font=\scriptsize] {sideways};
    \draw[go, thick] (0,0) -- (0,22) node[above, font=\scriptsize] {forward};
    \draw[gray] (5,0) arc (0:90:5);
    \node[font=\scriptsize] at (9,7) {\SI{90}{\degree}};
  \end{scope}

  % measured
  \begin{scope}[shift={(62,0)}]
    \node[font=\small\bfseries, anchor=south] at (18,26) {Where the arm was sent};
    \draw[go, thick, red!65!black] (0,0) -- (24,0)
      node[right, font=\scriptsize] {first movement};
    \draw[go, thick, red!65!black] (0,0) -- (-16.3,14.8)
      node[above left, font=\scriptsize, align=center] {second\\ movement};
    \draw[gray] (6,0) arc (0:137.7:6);
    \node[font=\scriptsize, red!65!black] at (7,11) {\SI{137.7}{\degree}};
  \end{scope}

  \node[tag, text width=112mm, align=left] at (40,-14)
    {Two directions that ought to be at right angles arrive
     \SI{47.7}{\degree} apart from where they should be. This is a property of
     the geometry, not of the labelling: three directions that are not
     mutually perpendicular cannot be made perpendicular by swapping or
     negating axes, and no setting in the software can do it. Fixing it means
     moving a sensor or adding one.};

\end{tikzpicture}
